\section{Transport absorption meets an adaptive index carrier}
The following is a verified incremental composition, not a qualitative
breakthrough.  It gives thirteen positive variables with degree $30$ and
$94$ arithmetic operations, and twelve positive variables with degree $114$
and $105$ operations.  The respective incoming points are $34/93$ and
$118/104$.  The operation increase is explicit.

Use the complete compact source44/source58 coefficient tuple, with
$\beta=2^{590069}-1$.  In the incoming congruence-carrier construction replace
only its second carrier by
\[
 K=4qY,\quad U=qK,\quad
 D_*=U(X+ac)+H(U\gamma+KX-r_0),\qquad r_0=k-R-1.
\]
The first transport carrier, the normalized strong equation, and the auxiliary
norm remain unchanged.  In particular the thirteen-variable residuals are
$A_0,A_1,A_a,A_s$, with $A_1=D_*^2-U^2(1+\Delta c^2)$, and the equation is
$A_0^2+A_1^2+A_a^2+A_s^2=0$.  Their exact degrees are $15,13,13,14$.
The supplied variables and all polynomial registers are listed in the literal
JSON graphs; no register introduces an additional witness.

\paragraph{The order of recovery.}
The unchanged first carrier recovers $q^3\mid Y$, positive transport, and
$C>0$, before using $q\mid X$.  Thus $3q<R<q^4-2$ and
$\gcd(q,H)=1$.  The new second carrier is positive, and integrality of its
square root gives $U\mid H(KX-r_0)$.  Since
\[
 d=\gcd(K,H)\le3,\qquad
 \mu=K/d\ge4q^4/3>R+1,
\]
we obtain $\mu\mid r_0$.  The normalized strong equation and the unchanged
signed auxiliary-index argument give $p=R$ without the syntactic shift
$X=x+q$ needed by the different ordinary-strong endpoint in the incoming
SharedTaps packet.  The first Pell congruence then gives
$\mu\mid2n-R-1$.  As $1\le n<R$, this forces $2n=R+1$ exactly.
Now $r_0$ is a positive multiple of $4XY^2$, and $U\mid4XY^2$ because
$q^3\mid Y$.  Returning to the full carrier yields $q\mid HX$, hence
$q\mid X$.  The recovered main-root coordinate
$\gamma_{\rm old}=\gamma+X/q-r_0/U$ is positive because
$D>(a+1)c$ and $c>X$.  All native source coordinates are restored.

\paragraph{Completeness and the twelve-variable endpoint.}
For a canonical source witness take
$\gamma=\gamma_{\rm old}-X/q+r_0/U>0$; the inherited canonical margin
$\gamma_{\rm old}>X/q$ suffices.  Every other witness is unchanged.
The incoming divided auxiliary norm still replaces its two coordinates by
one and yields the twelve-positive equation
$(1+A_0)(1+A_1)(1+N)-1=0$.  Its residual degrees are $15,13,86$.
Because $q$ and $U$ are even, all three residuals are squares modulo four;
therefore the unit-product finalizer forces all three residuals to vanish.
The exact degree is $114$.  The corresponding sum-of-squares sibling has
unchanged degree $172$ and is not presented as a new degree improvement.
The ordinary mathematical reduction, with the explicitly inherited compiler
and calibration, proves in ZFC that absence of positive zeros is equivalent
to $\operatorname{Con}(\mathrm{ZFC})$.

\paragraph{A failed additional deletion worth recording.}
Merging the two outer roots through a rational norm of $\lambda\tau+D$
would remove another witness if it preserved all residue channels.
A weight $\lambda=H(q-1)+1$ recovers transport but loses
$D\equiv X+ac\pmod H$.  A weight divisible by $H$ retains that residue but
only gives $(q-1)\mid H B_t$, not $(q-1)\mid B_t$.
These are inequivalent: $q=32$, $Y=q^3$, and $X=192$ give
$31\mid H=4Y(X+1)+3$.  Imposing $(q-1)\mid Y$ would repair coprimality but
has no justified canonical-completeness proof.  Accordingly no eleven-variable
endpoint is claimed from this idea.  The obstruction is recorded to prevent a
superficially attractive but invalid root merger.

The supplied checks verify exact local polynomial identities, degree,
liveness, and operation counts.  They are not consistency tests and do not
replace unbounded proofs or constitute a new proof-assistant certificate.
